%!TEX TS-program = xelatex
%!TEX encoding = UTF-8 Unicode
% Awesome CV LaTeX Template for Cover Letter
%
% This template has been downloaded from:
% https://github.com/posquit0/Awesome-CV
%
% Authors:
% Claud D. Park <posquit0.bj@gmail.com>
% Lars Richter <mail@ayeks.de>
%
% Template license:
% CC BY-SA 4.0 (https://creativecommons.org/licenses/by-sa/4.0/)
%


%-------------------------------------------------------------------------------
% CONFIGURATIONS
%-------------------------------------------------------------------------------
% A4 paper size by default, use 'letterpaper' for US letter
\documentclass[11pt, a4paper]{awesome-cv}

% Configure page margins with geometry
\geometry{left=1.4cm, top=.8cm, right=1.4cm, bottom=1.8cm, footskip=.5cm}

% Specify the location of the included fonts
\fontdir[fonts/]

% Color for highlights
% Awesome Colors: awesome-emerald, awesome-skyblue, awesome-red, awesome-pink, awesome-orange
%                 awesome-nephritis, awesome-concrete, awesome-darknight
\colorlet{awesome}{awesome-skyblue}
% Uncomment if you would like to specify your own color
% \definecolor{awesome}{HTML}{CA63A8}

% Colors for text
% Uncomment if you would like to specify your own color
% \definecolor{darktext}{HTML}{414141}
% \definecolor{text}{HTML}{333333}
% \definecolor{graytext}{HTML}{5D5D5D}
% \definecolor{lighttext}{HTML}{999999}

% Set false if you don't want to highlight section with awesome color
\setbool{acvSectionColorHighlight}{true}

% If you would like to change the social information separator from a pipe (|) to something else
\renewcommand{\acvHeaderSocialSep}{\quad\textbar\quad}


%-------------------------------------------------------------------------------
%	PERSONAL INFORMATION
%	Comment any of the lines below if they are not required
%-------------------------------------------------------------------------------
% Available options: circle|rectangle,edge/noedge,left/right
\photo[circle,noedge,left]{./profile_Bia.png}
\name{Molnár-Bebtó}{Bianka}
\position{Pszichológus hallgató{}}
\address{2132 Göd Család utca 22/B.}

\mobile{(+36) 30-444-98-48}
\email{b.bebto@gmail.com}
% \homepage{www.posquit0.com}
% \github{posquit0}
% \linkedin{posquit0}
% \gitlab{gitlab-id}
% \stackoverflow{SO-id}{SO-name}
% \twitter{@twit}
% \skype{skype-id}
% \reddit{reddit-id}
% \medium{madium-id}
% \googlescholar{googlescholar-id}{name-to-display}
%% \firstname and \lastname will be used
% \googlescholar{googlescholar-id}{}
% \extrainfo{extra informations}

% \quote{``Be the change that you want to see in the world."}


%-------------------------------------------------------------------------------
%	LETTER INFORMATION
%	All of the below lines must be filled out
%-------------------------------------------------------------------------------
% The company being applied to
\recipient
  {Szegedi Tudományegyetem Bölcsészet- és Társadalomtudományi Kar}
  {\\\\}
% The date on the letter, default is the date of compilation
\letterdate{\selectlanguage{magyar}\today}
% The title of the letter
\lettertitle{Motivációs levél a pszichológia mesterszak felvételi eljáráshoz}
% How the letter is opened
% \letteropening{Kedves Zsuzsanna}
% How the letter is closed
\letterclosing{Üdvözlettel,}
% Any enclosures with the letter
\letterenclosure[Csatolva]{Curriculum Vitae}


%-------------------------------------------------------------------------------
\begin{document}

% Print the header with above personal informations
% Give optional argument to change alignment(C: center, L: left, R: right)
\makecvheader[R]

% Print the footer with 3 arguments(<left>, <center>, <right>)
% Leave any of these blank if they are not needed
\makecvfooter
  {\selectlanguage{magyar}\today}
  {Molnár-Bebtó Bianka~~~·~~~Motivációs levél}
  {}

% Print the title with above letter informations
\makelettertitle

%-------------------------------------------------------------------------------
%	LETTER CONTENT
%-------------------------------------------------------------------------------
\begin{cvletter}

\lettersection{Rólam}
Ezúton szeretnék felvételt nyerni a SZTE-BTK Pszichológia szakának a mesterképzési szintjére. Jelenleg a Károli Gáspár Református Egyetem Bölcsészet- és Társadalom Tudományi Kar Pszichológia szakának alapkézésén vagyok végzős hallgató. Régi álmom vált valóra, amikor felvételt nyertem és elkezdhettem pszichológiai tanulmányaimat. Az elmúlt három év alatt folyamatosan erősödött bennem a hivatástudat. A tanárokkal, szakemberekkel és diáktársaimmal való találkozás, megismerkedés nagyon inspirálóan és motiválóan hatott és hat rám a mai napig. Nagyon örülök, hogy ennek a világnak a részese lehetek, melynek alapja a folyamatos fejlődés, önképzés és önismeret. Ami alkalmassá tesz a segítő pályára, a belső motivációm, érdeklődésem, elhivatottságom és nyitott, empatikus személyiségem.


\lettersection{Miért a SZTE-BTK?}
Bízom abban, hogy az Önök egyeteménél szerzett tudás és tapasztalat nagyban hozzájárul ahhoz, hogy jó segítőszakember váljon belőlem. Szeretném látásmódomat bővíteni azáltal, hogy a mesterképzést már Önöknél végezném el. Hiszem, hogy még több tapasztalatot és tudást szerezhetek más tanárokkal és szakemberekkel való közös munka során.


\lettersection{Melyik specializáció?}
Elsősorban a Tanácsadás és iskolapszichológia érdekel, de a klinikai vonal is nagyon foglalkoztat. A klinikai specializáción belül a gyermek és a felnőtt között nem tudok különbséget tenni, mert mindkettőt ugyanolyan izgalmasnak találom.

\end{cvletter}


%-------------------------------------------------------------------------------
% Print the signature and enclosures with above letter informations
\makeletterclosing

\end{document}
